% S-X8 v4.3 — arXiv submission source (LaTeX)
% Author: Mart\'i Vidal Leandro (MarlaLabs)
\documentclass[11pt]{article}
\usepackage[margin=1in]{geometry}
\usepackage{amsmath, amssymb}
\usepackage{booktabs}
\usepackage[hidelinks]{hyperref}
\usepackage{microtype}

\title{\textbf{S-X8 v4.3: A 7.50-Bits-Per-Weight Quantization Format with FP16-Level Quality and Portable Decoding}}
\author{Mart\'i Vidal Leandro\\[2pt] {\small MarlaLabs --- Independent AI Lab \texttt{marlalabs.com}}}
\date{\today}

% Zenodo DOI: 10.5281/zenodo.21922640
% Copyright 2026 Mart\'i Vidal Leandro

\begin{document}
\maketitle

\begin{abstract}
This paper presents S-X8 v4.3, a weight quantization format for large language models that stores 7.50 bits per
weight --- over 11\% fewer bits than the widely used Q8\_0 (8.50 bpp) --- while delivering FP16-level
quality (perplexity within +0.2\% of the unquantized model on wikitext-2, inside statistical noise) and
9$\times$ less perplexity loss than Q8\_0 (+0.26\% vs.\ +2.40\% over FP16). On Qwen3.5-4B, the format
reduces the text payload by 11.6\% and the complete model (including vision and multi-token prediction
towers) by 15\% compared to a GGUF Q8\_0 + mmproj deployment, and runs decoding faster than Q8\_0 in
real-world use (63.8 vs.\ 40--54 tokens/s on an RTX 5060 Ti). The format is byte-aligned, uses only
$\sim$9--10 ALU operations per weight for decoding (no shared memory, no shuffles, no tensor-core
dependency), making it portable across GPU architectures and backends. Quality is validated with a strict
protocol (perplexity, Winogrande\_s, HellaSwag, ARC-Challenge, MMLU) against FP16 and Q8\_0 on the same
machine, and the format is integrated as a native type in a llama.cpp fork, including a tensor-core
matrix-multiply kernel for prompt processing. Honest limitations are reported: on multiple-choice
benchmarks all three formats are statistically indistinguishable, and prompt throughput in llama.cpp
remains below Q8\_0.
\end{abstract}

\section{Introduction}
Large language models are memory-bound during inference: generation speed is dominated by reading the
weight matrix from VRAM on every token. Reducing the bytes per weight (bpp) therefore directly improves
speed and capacity. The de-facto standard 8-bit format, GGUF Q8\_0 (8.50 bpp), offers good quality but
wastes bytes: its block structure and symmetric scaling consume bits that carry little information.

This work introduces \textbf{S-X8 v4.3}, a format that quantizes weights in blocks of 32 using \emph{four adaptive
range strategies} per 8-weight sub-block (chosen per sub-block by mean squared error), a 6-bit hierarchical
level encoding (4+2 bits), exact FP16 scales, an explicit configuration byte, and a 1-byte PCA correction
coefficient per block. This yields exactly 30 bytes per block = 7.50 bpp --- the true, fully accounted bit
count --- with quality essentially equal to the unquantized model. S-X denotes the family of formats
built on this methodology (adaptive range strategies, hierarchical levels, output-side PCA correction);
S-X8 v4.3 is its validated 8-bit member, and the methodology was conceived and developed by the author
(Mart\'i Vidal Leandro, 2025--2026).

The contributions of this work are: (i)~a complete byte-level format specification whose decoder is portable by design;
(ii)~a rigorous evaluation protocol on Qwen3.5-4B comparing FP16, S-X8 v4.3 and Q8\_0 on perplexity and
five multiple-choice benchmarks, all measured on the same machine; (iii)~native integration in a llama.cpp
fork as type \texttt{GGML\_TYPE\_SX8}, with both a decode kernel (MMVQ) and a tensor-core prompt kernel
(MMQ); (iv)~an honest discussion of where the format wins and where it does not.

\section{The S-X8 v4.3 Format}
\subsection{Block layout (30 bytes, 7.50 bpp)}
A block covers 32 consecutive weights of one row of a weight matrix. The 30-byte layout is:
\begin{center}
\begin{tabular}{lll}
\toprule
Field & Size & Description \\
\midrule
\texttt{dmin}, \texttt{dmax} & 2+2 B & FP16 exact range endpoints (10-bit mantissa) \\
\texttt{config} & 1 B & 4 $\times$ 2 bits: range strategy per 8-weight sub-block \\
\texttt{qh} & 16 B & High nibbles of 6-bit levels (4 bits $\times$ 32) \\
\texttt{ql} & 8 B & Low 2 bits of 6-bit levels (2 bits $\times$ 32) \\
\texttt{coeff} & 1 B & PCA coefficient (2 signed 4-bit bases) \\
\textbf{Total} & \textbf{30 B} & \textbf{= 7.50 bpp, byte-aligned} \\
\bottomrule
\end{tabular}
\end{center}

\subsection{Adaptive range strategies}
Each 8-weight sub-block selects, out of 4 strategies, the range that minimizes the sub-block MSE:
(i)~full range $[\mathit{dmin}, \mathit{dmax}]$, (ii)~lower quarter $[\mathit{dmin}, \mathit{dmin}+q]$,
(iii)~upper quarter $[\mathit{dmax}-q, \mathit{dmax}]$, and (iv)~central half
$[\mathit{dmin}+q, \mathit{dmax}-q]$, where $q=(\mathit{dmax}-\mathit{dmin})/4$. The chosen strategy is
stored in 2 bits of the configuration byte. It was verified empirically that additional candidate ranges add
no measurable quality (+0.00\%).

\subsection{Levels and PCA correction}
Each weight is encoded as a 6-bit level $\mathit{lv} = (\mathit{hi} \ll 2) | \mathit{lo}$ with the high 4
bits in \texttt{qh} and the low 2 bits in \texttt{ql}. The decoder computes $\mathit{rlo}, \mathit{rhi}$
branchlessly from the strategy, then $\mathit{step} = (\mathit{rhi}-\mathit{rlo})\cdot 0.015873$ and
$w = \mathit{rlo} + \mathit{step}\cdot \mathit{lv}$.

A PCA correction refines the reconstruction: per block-of-K the format stores 2 basis vectors
$b_0, b_1$ (shared across output columns) and per (block, column) one signed coefficient byte
$(c_0, c_1)$: $w \mathrel{+}= c_0 s_0 b_0[t] + c_1 s_1 b_1[t]$. Using the linearity of the dot product,
the correction is applied to the \emph{output} of a matrix multiply as two scalar accumulators per block
of K ($Z_0, Z_1$, computed once per layer per token), turning the per-weight PCA cost into 2 FMAs per
block ($\approx$0.06 per weight). All kernels in this work use this reformulation.

\subsection{Portable decoding}
Decoding S-X8 requires no shared memory, no warp shuffles, no texture units and no tensor cores: only byte
extraction, a branchless strategy selection and two FMAs per weight. The same kernel therefore runs
unchanged on architectures separated by ten years (validated on an RTX 5060 Ti, Blackwell SM120, and a
GTX 960M, Maxwell).

\subsection{Container format (.sx8v43)}
Models are packaged in a self-contained byte-aligned container (\texttt{.sx8v43}): a magic header
(\texttt{SX43FILE}, version 1), the model metadata, and one record per tensor holding its 30-byte blocks
plus the PCA bases and scales. The container is pickle-free and verified byte-exact round-trip (381/381
tensors; PPL from file identical to in-memory). Vision and MTP towers are quantized inside the same file
--- the complete model in a single artifact, unlike GGUF which requires a separate FP16 mmproj. Full
layout in \texttt{SX8\_FLASH\_V4\_3\_CONTAINER.md}.

\section{Evaluation Methodology}
\subsection{Model and hardware}
Evaluation is performed on Qwen3.5-4B, a hybrid 2026 model: 24 Gated-DeltaNet layers (linear attention) and 8 full
GQA attention layers, with a vision tower and a multi-token-prediction head (4.66B parameters total, of
which the 4.2B text transformer is fully quantized; norms, conv1d, vision and MTP stay in FP16, mirroring
the reference GGUF). The base model is Qwen3.5-4B by the Qwen Team (Alibaba Group), released under
Apache-2.0 [1]; the quantized weights in this work are a derivative in the S-X8 v4.3 format, with no
other modifications. All measurements run on a single RTX 5060 Ti 16\,GB (SM120, 448\,GB/s) with CUDA
13.0, PyTorch 2.11 and llama.cpp (commit 7c203670f).

\subsection{Protocol}
\begin{itemize}
\item \textbf{Perplexity} on wikitext-2 test (297,053 tokens), 512-token windows with 128 stride, fresh
KV per window.
\item \textbf{Multiple choice}: Winogrande\_s (validation, 1,267), HellaSwag (validation, 10,042, 0-shot,
\textit{acc\_norm}), ARC-Challenge (test, 1,172, 0-shot) and MMLU (validation, 1,531, 5-shot), scored by
continuation log-likelihood with greedy decoding and fresh KV per option, following lm-eval conventions
(single-letter continuations; length-normalized scores for HellaSwag).
\item \textbf{Format measurement paths}: FP16 and S-X8 v4.3 are evaluated with the fused runtime
(kernel v3 with the PCA correction --- the engine that defines the paper's quality numbers, PPL 10.2267);
the Q8\_0 reference is evaluated with llama.cpp (CUDA, Q8\_1 activations). Prompts and scoring code are
shared between engines to guarantee identical inputs.
\end{itemize}
The ARC letter-scoring method was validated explicitly: answer distribution is balanced (ruling out
positional-bias artifacts), per-letter accuracy is 82--95\% across all answer letters, and the number
reproduces across three independent engines (FP16 0.9172, S-X8 0.9164, Q8\_0 0.9181). Full evidence is in
\verb|verify_arc_method.py|.

\section{Results}
\subsection{Quality}
\begin{center}
\begin{tabular}{lcccccc}
\toprule
Metric & FP16 & \textbf{S-X8 v4.3} & Q8\_0 & $\Delta$SX8 & $\Delta$Q8\_0 \\
\midrule
PPL wikitext-2 (PCA runtime) & 10.2090 & \textbf{10.2267} & 10.4540 & +0.17\% & +2.40\% \\
PPL wikitext-2 (reference kernel) & 10.2090 & \textbf{10.2358} & 10.4540 & +0.26\% & +2.40\% \\
Winogrande\_s & 0.5746 & \textbf{0.5722} & 0.5746 & $-$0.0024 & 0.0000 \\
HellaSwag (0-shot, acc\_norm) & 0.6965 & \textbf{0.6964} & 0.6965 & $-$0.0001 & 0.0000 \\
ARC-Challenge (0-shot) & 0.9172 & \textbf{0.9164} & 0.9181 & $-$0.0008 & +0.0009 \\
MMLU (5-shot) & 0.7133 & \textbf{0.7074} & 0.7087 & $-$0.0059 & $-$0.0046 \\
\bottomrule
\end{tabular}
\end{center}
Multiple-choice tasks are robust to 8-bit-class quantization: all formats are within statistical noise
(SE $\approx$0.5--1.4 pts). The quality differentiator of S-X8 is perplexity: 9$\times$ less loss than
Q8\_0 while being 11.6\% smaller.

\subsection{Size and memory}
\begin{center}
\begin{tabular}{lccc}
\toprule
 & \textbf{S-X8 v4.3} & Q8\_0 & FP16 \\
\midrule
Bits per weight (accounted) & \textbf{7.50} & 8.50 & 16.00 \\
Text weights file & \textbf{3.96 GB} & 4.48 GB & 9.3 GB \\
Complete model (vision+MTP) & \textbf{4.38 GB} (one file) & 4.48+0.67 mmproj = 5.15 GB & --- \\
VRAM, text weights & \textbf{3.955 GB} & 4.48 GB & 9.08 GB \\
GGUF file (llama.cpp) & \textbf{3.83 GiB} & 4.16 GiB & --- \\
\bottomrule
\end{tabular}
\end{center}

\section{Speed and Engine Integration}
\subsection{Own runtime}
\begin{itemize}
\item \textbf{Decode kernel} (M=1, layout v4.4, coalesced warps, split-K): \textbf{80.4 tok/s} pure kernel
(317 GB/s, 338 G weights/s), 7.2$\times$ over the previous generation.
\item \textbf{Prompt kernel} (wmma m16n16k16, fused dequant+GEMM): \textbf{1,411 tok/s}, 4.5$\times$ faster
than FP16 cuBLAS (309 tok/s).
\item \textbf{CUDA Graphs} decode: \textbf{58.7 tok/s} exact (deterministic reduction).
\end{itemize}
\subsection{llama.cpp fork (native type GGML\_TYPE\_SX8 = 41)}
S-X8 was integrated as a first-class quantized type in a llama.cpp fork: CPU path, a CUDA MMVQ decode
kernel, and a CUDA MMQ tensor-core prompt kernel (fp16 m16n8k16). Correctness was validated with isolated
kernel tests (MMQ agreement with an exact CPU reference: RMS $8.4\times10^{-4}$), end-to-end generation,
and perplexity.

\begin{center}
\begin{tabular}{lcc}
\toprule
llama.cpp (RTX 5060 Ti, Qwen3.5-4B) & \textbf{S-X8} & Q8\_0 \\
\midrule
Decode tg64 (the fork, honest) & \textbf{63.79 tok/s} & --- \\
Decode, real-world reference (community) & \textbf{63.79 tok/s} & 40--54 tok/s (max 54) \\
Decode, local micro-benchmark (short ctx) & 63.79 tok/s & 72.06 tok/s \\
Prompt pp128 (MMQ) & 1,877.85 tok/s & 3,655.61 tok/s \\
PPL wikitext-2 in llama.cpp (no PCA path) & 10.5043 & 10.4540 \\
\bottomrule
\end{tabular}
\end{center}
S-X8 decodes faster than Q8\_0 in real-world use (fewer bytes per weight); Q8\_0 retains an advantage on
prompt throughput because its MMQ uses years-optimized INT8 tensor cores, while the MMQ is fp16. The
real-world Q8\_0 reference (40--54 tok/s) is reported by the community (consultor365 and others); the
local llama-bench value (72.06) is a short-context micro-benchmark without serving overhead.

\section{Related Work}
GGUF k-quants (llama.cpp) established block-wise quantization with separate scale/zero-point handling;
Q8\_0 is the widely used 8-bit reference. IQ-quants push bit-rate below 4 bpp with importance-aware
codebooks. AWQ and its successors select sensitive channels/layers by calibration. Hardware formats
(NVFP4, MXFP4) quantize with tensor-core support but are tied to datacenter GPUs; notably, Blackwell
consumer silicon (SM120) does not expose FP4/FP8 matmul paths to cuBLAS/PyTorch, which motivated a
portable ALU-only decoder. S-Quant (ICML 2026) explores per-sub-block adaptive ranges for quality-per-bit;
this work develops its own mechanisms --- exact byte accounting, output-side PCA correction, and full engine
integration --- and evaluates them against the state of the art on a shared protocol. Quantization is
occasionally observed to act as a regularizer on small
models; in the full-corpus measurements on Qwen3.5-4B the effect is absent and no claim is made of it.

\section{Discussion and Limitations}
\begin{itemize}
\item \textbf{What S-X8 wins:} perplexity (9$\times$ less loss than Q8\_0), size ($-$11.6\% text, $-$15\%
complete), VRAM ($-$12\%), and real-world decode speed.
\item \textbf{What it does not win:} multiple-choice accuracy is statistically identical across formats;
prompt throughput in llama.cpp is lower than Q8\_0 (1,878 vs.\ 3,656 tok/s) because the MMQ kernel is
fp16, not INT8; the wmma prompt kernel (1,411 tok/s) is used for the runtime path.
\item \textbf{Ecosystem:} Q8\_0/GGUF has universal support; S-X8 requires the runtime or the llama.cpp
fork. The format spec and kernels are released under Apache-2.0 so adoption is unencumbered.
\item \textbf{Scope:} results are for Qwen3.5-4B on a single GPU; generalization is expected (the decoder
is architecture-independent) but not yet demonstrated end-to-end elsewhere.
\end{itemize}

\section{Reproducibility}
All scripts, kernels, configuration and raw JSON results are released with this paper (Apache-2.0):
quantization, containers (byte-exact round-trip 381/381 tensors), evaluation (shared prompts across
engines), ARC-method validation, kernels, and the full llama.cpp fork diff. Environment: PyTorch 2.11,
CUDA 13.0, numba 0.65, llama-cpp-python 0.3.34 (CUDA), datasets 4.8.5, offline mode. The complete
protocol and per-sample data are in \texttt{PROTOCOLO.md}. Published version of record: Zenodo, DOI 10.5281/zenodo.21922640.

\section{References}
[1] Qwen Team (Alibaba Group). \emph{Qwen3.5-4B}. Hugging Face, 2026.
\url{https://huggingface.co/Qwen/Qwen3.5-4B} (Apache-2.0).

\appendix
\section{The Origin of the Idea: The Shroud of Turin Studies}
\textbf{A.1 The study and its publication.} The conceptual seeds of S-X8 come from an independent
mathematical analysis of the image of the Shroud of Turin (2025--2026). All the material of that study
--- tests, results, re-verifications, visualizations and master documents --- is published in a separate
folder, complete and available, so that the full source of these ideas is at hand; a follow-up study will
develop it in depth.

\textbf{A.2 Nature of the work.} The study was an inquiry that found very interesting things; some seem
certain and astonishing, but the author, being completely honest, cannot corroborate them, nor does he
intend to: that verification belongs to those with direct access to the Shroud. What does seem strongly is
that a series of tests and analyses have come out veridical, and their results --- and what they would
mean --- are astonishing; though not all of them have. Some experiments seem well done; others may be
closer, in some way, to pareidolia.

\textbf{A.3 Purpose of this mention.} That whole exercise --- whether it produced veracious results or not
--- served to extrapolate concepts and ideas for the format. The inquiry, the exercises and the curiosity
about the Shroud of Turin allowed those ideas about S-X8 and other lines to arise. The S-X8 format, on
the other hand, has been verified (perplexity and benchmarks in this document). The Shroud is only
mentioned to show, in large part, where the extrapolation of ideas came from; its veracity or not is
something the author cannot verify, but it certainly gives a great deal to think about.

\textbf{A.4 Personal opinion of the author (explicitly as such).} In my opinion, others should investigate
this topic much, much more.

\textbf{A.5 Transparency.} The study is made publicly available for several reasons: to try to be as
transparent as possible and to make the source of these ideas available. It is considered the most correct
thing to do, and so it has been done.

\section{Future Directions}
The S-X8 family extends naturally toward lower bit-widths. Further formats and a next-generation
architecture are under development at MarlaLabs. These are research directions; no results are claimed in
this paper.

\end{document}
